\documentclass[10pt,a4paper]{amsart}
\usepackage[margin=19mm]{geometry}
\usepackage{amsmath,amssymb,mathtools}
\usepackage[scr=rsfs,cal=euler]{mathalfa}
\usepackage{xfrac}
\usepackage[unicode,hidelinks,bookmarks=false]{hyperref}
\newcommand{\ie}{i.e.\ }
\newcommand{\cf}{cf.\ }
\newcommand{\resp}{respectively\ }
\newcommand{\Subsection}{\subsection}
\providecommand{\nobreakdash}{\nobreak\hbox{-}}
\setlength{\emergencystretch}{2em}
\multlinegap0pt
\usepackage{amssymb}
\usepackage{mathtools}
\usepackage[shortlabels]{enumitem}
\newtheorem{lemma}{Lemma}
\title{The Irrationality Measure of $\pi$ Is at Most $7.101862832357$}
\author{Yufei Bai}
\date{Source: arXiv:2609.11276v2 (CC BY 4.0)}
\begin{document}
\maketitle

\noindent\emph{Source and licence.} The Irrationality Measure of $\pi$ Is at Most $7.101862832357$. Version arXiv:2609.11276v2 (CC BY 4.0), distributed by arXiv under the Creative Commons Attribution 4.0 International licence, http://creativecommons.org/licenses/by/4.0/, as recorded in the arXiv licence metadata for this exact version and shown on the abstract page https://arxiv.org/abs/2609.11276v2. Full licence notice: \texttt{licenses/arxiv-2609.11276-CC-BY-4.0.txt}.\par\smallskip

\noindent\textit{Editorial scope.} Counted unit: the article's lemma lem:hata-one-number (source0.tex lines 1596--1612). The article's complete original proof of it is delivered with it (source0.tex lines 1614--1748, one \texttt{proof} environment of 135 lines). No mathematics is written, repaired or re-derived: the statement and the proof are the article's own lines, in the article's own notation, and no retained line is substituted or re-flowed.  No further source-local context is needed: every cross-reference of every retained line resolves inside the two delivered spans.  Nothing was omitted from any retained span, so the package carries no omission record.  No retained line contains a citation, so the wrapper carries no bibliography and no bibliography entry.  No retained line contains a figure environment, an image-inclusion command or drawing code, so no image is delivered and the package declares no asset dependency.  Editorial formatting only, in addition: an AMS \texttt{amsart} preamble that restates the article's own declarations and macros used by the retained lines, namely the environments \texttt{lemma} on separate counters, together with the standard packages those lines need.  The retained environments print in this document as Lemma 1 in order of appearance, whereas the article prints the same items with the numbering of its own full document.  The complete native source of the article as distributed in the arXiv e-print for this exact version is bundled unchanged as \texttt{sources/209998/source0.tex}.
\begin{lemma}[Hata's index-selection lemma]\label{lem:hata-one-number}
Let \(\theta\) be irrational.  Suppose that \(U_n,V_n\in\mathbb Z\),
that \(V_n\ne0\) for every sufficiently large \(n\), and put
\[
 \Lambda_n=U_n+V_n\theta.
\]
If there are \(\sigma>0\) and \(\tau>0\) such that
\begin{equation}\label{eq:hata-hypotheses}
 \lim_{n\to\infty}\frac1n\log|V_n|=\sigma,
 \qquad
 \limsup_{n\to\infty}\frac1n\log|\Lambda_n|\leq-\tau,
\end{equation}
then
\begin{equation}\label{eq:hata-conclusion}
 \mu(\theta)\leq1+\frac{\sigma}{\tau}.
\end{equation}
\end{lemma}

\begin{proof}
Because \(V_n\ne0\) eventually and \(\theta\notin\mathbb Q\), the form
\(\Lambda_n\) is also nonzero at every sufficiently large index.  Thus
the logarithms in \eqref{eq:hata-hypotheses} are defined after finitely
many terms have been discarded.

Fix \(\varepsilon>0\).  Choose \(\delta\) so that
\begin{equation}\label{eq:hata-delta-choice}
 0<\delta<\min\left\{\frac\sigma2,\frac\tau6\right\},
 \qquad
 \kappa:=\frac{\sigma+\delta}{\tau-3\delta}
 <\frac\sigma\tau+\frac\varepsilon2.
\end{equation}
This is possible by continuity at \(\delta=0\).  The ordinary limit and
the limsup in \eqref{eq:hata-hypotheses} give an index \(n_0\) such that,
for every \(n\geq n_0\),
\begin{equation}\label{eq:hata-exponential-bounds}
 e^{(\sigma-\delta)n}\leq|V_n|
 \leq e^{(\sigma+\delta)n},
 \qquad
 |\Lambda_n|\leq e^{-(\tau-\delta)n}.
\end{equation}

Let \(p\in\mathbb Z\) and \(q\geq1\), and write
\begin{equation}\label{eq:hata-determinant-identity}
 \Delta=q\theta-p\ne0,
 \qquad A_n=qU_n+pV_n\in\mathbb Z,
 \qquad q\Lambda_n=A_n+V_n\Delta.
\end{equation}
There are infinitely many indices \(n\geq n_0\) for which \(A_n\ne0\).
Indeed, if \(A_n=0\) eventually, then \eqref{eq:hata-determinant-identity}
and \eqref{eq:hata-exponential-bounds} would imply
\[
 0<|\Delta|=q\frac{|\Lambda_n|}{|V_n|}
 \leq q e^{-(\sigma+\tau-2\delta)n}\longrightarrow0,
\]
which is impossible.

Select the first index at the scale determined by \(q\):
\begin{equation}\label{eq:hata-first-index}
 N=N(q):=\min\{n\geq1:2q<e^{(\tau-\delta)n}\}.
\end{equation}
For all sufficiently large \(q\), one has \(N\geq n_0\).  Minimality
of \(N\) gives
\begin{equation}\label{eq:hata-first-index-size}
 e^{(\tau-\delta)(N-1)}\leq2q<e^{(\tau-\delta)N},
 \qquad
 e^N\leq e(2q)^{1/(\tau-\delta)}.
\end{equation}
Moreover, for every \(n\geq N\),
\begin{equation}\label{eq:hata-small-form-at-scale}
 q|\Lambda_n|\leq q e^{-(\tau-\delta)n}<\frac12.
\end{equation}

Suppose first that \(A_N\ne0\).  Since it is an integer,
\eqref{eq:hata-determinant-identity} and
\eqref{eq:hata-small-form-at-scale} yield
\[
 |V_N|\,|\Delta|=|q\Lambda_N-A_N|>\frac12.
\]
Using \eqref{eq:hata-exponential-bounds} and
\eqref{eq:hata-first-index-size}, we obtain
\begin{equation}\label{eq:hata-case-one}
 |\Delta|>\frac12e^{-(\sigma+\delta)N}
 \geq c_1q^{-\kappa_1},
 \quad
 \kappa_1=\frac{\sigma+\delta}{\tau-\delta}<\kappa,
 \quad
 c_1=2^{-1-\kappa_1}e^{-(\sigma+\delta)}>0.
\end{equation}

It remains to handle the case \(A_N=0\).  By the infinitude just proved,
there is a first index
\[
 M=\min\{m>N:A_m\ne0\}.
\]
Then \(A_{M-1}=0\).  At this preceding index,
\eqref{eq:hata-determinant-identity} and both coefficient estimates in
\eqref{eq:hata-exponential-bounds} give
\begin{equation}\label{eq:hata-zero-predecessor}
 |\Delta|=q\frac{|\Lambda_{M-1}|}{|V_{M-1}|}
 \leq q e^{-(\sigma+\tau-2\delta)(M-1)}.
\end{equation}
Set
\[
 \gamma=\sigma+\tau-2\delta,
 \qquad \rho=\frac{\sigma+\delta}{\gamma}.
\]
The choice \(\delta<\tau/6\) gives
\(\gamma-(\sigma+\delta)=\tau-3\delta>0\), and hence
\(0<\rho<1\).  Rearranging \eqref{eq:hata-zero-predecessor} gives
\begin{equation}\label{eq:hata-second-index-size}
 e^M\leq e\left(\frac q{|\Delta|}\right)^{1/\gamma}.
\end{equation}
At the index \(M\), the integer \(A_M\) is nonzero, so the same argument
as above, now combined with \eqref{eq:hata-second-index-size}, yields
\[
 |\Delta|>\frac12e^{-(\sigma+\delta)M}
 \geq\frac12e^{-(\sigma+\delta)}
 \left(\frac{|\Delta|}{q}\right)^\rho.
\]
Since \(1-\rho>0\), this is equivalent to
\begin{equation}\label{eq:hata-case-two}
 |\Delta|>c_2q^{-\rho/(1-\rho)}=c_2q^{-\kappa},
 \qquad
 c_2=\left(2e^{\sigma+\delta}\right)^{-\gamma/(\tau-3\delta)}>0.
\end{equation}
Here we used
\(\rho/(1-\rho)=(\sigma+\delta)/(\tau-3\delta)=\kappa\).
This second selection is essential: the zero at \(N\) controls the
distance to the next nonzero integer through \(|\Delta|\) itself.

Put \(c=\min\{1,c_1,c_2\}\).  For definiteness, take
\begin{equation}\label{eq:hata-uniform-threshold}
 Q(\varepsilon)=\left\lceil\max\left\{
 2,\ \frac12e^{(\tau-\delta)(n_0-1)},\ c^{-2/\varepsilon}
 \right\}\right\rceil.
\end{equation}
The middle term ensures \(N(q)\geq n_0\).  Thus both cases, uniformly
for every \(p\in\mathbb Z\) and every \(q\geq Q(\varepsilon)\), give
\[
 |q\theta-p|>cq^{-\kappa}
 >q^{-\sigma/\tau-\varepsilon};
\]
the last strict inequality follows from
\(\sigma/\tau+\varepsilon-\kappa>\varepsilon/2\) and the last term in
\eqref{eq:hata-uniform-threshold}.  Division by \(q\) proves
\[
 \left|\theta-\frac pq\right|
 >q^{-1-\sigma/\tau-\varepsilon}.
\]
All constants and the threshold are independent of \(p\).  Since
\(\varepsilon>0\) was arbitrary, the definition of \(\mu(\theta)\)
proves \eqref{eq:hata-conclusion}.
\end{proof}

\end{document}
